\chapter*{Resumo}

A identifica\c{c}\~ao e a organiza\c{c}\~ao espacial de parcelas experimentais em imagens orbitais e de drone constituem um desafio para a an\'alise automatizada de ensaios agr\'icolas, principalmente quando a densidade, a orienta\c{c}\~ao e a variabilidade visual das parcelas afetam o desempenho dos detectores. Esta pesquisa tem como objetivo investigar m\'etodos de vis\~ao computacional para detectar, localizar e organizar parcelas experimentais de forrageiras e avaliar como a composi\c{c}\~ao do conjunto de treinamento influencia a capacidade de generaliza\c{c}\~ao do detector. A metodologia compreende a constru\c{c}\~ao e a reorganiza\c{c}\~ao progressiva dos conjuntos de dados, o treinamento e a avalia\c{c}\~ao de modelos para segmenta\c{c}\~ao e detec\c{c}\~ao de parcelas, al\'em da investiga\c{c}\~ao de abordagens auxiliares para localiza\c{c}\~ao, classifica\c{c}\~ao da dificuldade das imagens e associa\c{c}\~ao espacial entre parcelas. As avalia\c{c}\~oes realizadas evidenciaram diferen\c{c}as de desempenho entre as representa\c{c}\~oes investigadas, com destaque das caixas delimitadoras orientadas em parte das m\'etricas, sem que essas diferen\c{c}as possam ser atribu\'idas exclusivamente ao tipo de representa\c{c}\~ao. Tamb\'em foi aplicado um procedimento de associa\c{c}\~ao geom\'etrica das parcelas entre imagens de um mesmo campo, que atribuiu identificador a 90,77\% dos objetos, sem que a corre\c{c}\~ao desses identificadores tenha sido medida; permanecem ainda limita\c{c}\~oes relacionadas \`{a} independ\^encia entre as imagens de treinamento e de teste. Como contribui\c{c}\~ao da pesquisa, ser\'a investigada uma estrat\'egia de sele\c{c}\~ao das imagens de treinamento orientada por um crit\'erio quantitativo de instabilidade das predi\c{c}\~oes, estimado em imagens n\~ao utilizadas no treinamento, e comparada ao conjunto completo, a subconjuntos aleat\'orios de mesmo tamanho e ao limiar hist\'orico de F1-Score de 0,97.

\textbf{Palavras-chave:} vis\~ao computacional; parcelas experimentais; fenotipagem vegetal; imagens orbitais; drones; detec\c{c}\~ao de objetos.

\chapter*{Abstract}

The identification and spatial organization of experimental plots in satellite and drone imagery constitute a challenge for the automated analysis of agricultural trials, particularly when plot density, orientation, and visual variability affect detector performance. This research aims to investigate computer vision methods for detecting, locating, and organizing experimental forage plots, and to assess how the composition of the training set influences the detector's generalization capability. The methodology comprises the construction and progressive reorganization of the datasets, the training and evaluation of models for plot segmentation and detection, as well as the investigation of auxiliary approaches for localization, image difficulty classification, and spatial association between plots. The evaluations carried out revealed performance differences among the representations investigated, with oriented bounding boxes standing out in some of the metrics, although these differences cannot be attributed exclusively to the type of representation. A procedure for the geometric association of plots across images of the same field was also applied, which assigned an identifier to 90.77\% of the objects, without the correctness of these identifiers having been measured; limitations remain regarding the independence between the training and test images. As a contribution of this research, a strategy for selecting training images guided by a quantitative criterion of prediction instability, estimated on images not used in training, will be investigated and compared against the full set, random subsets of the same size, and the historical F1-Score threshold of 0.97.

\textbf{Keywords:} computer vision; experimental plots; plant phenotyping; satellite imagery; drones; object detection.
